\section{Self-Training with Pseudo-Labels}
\label{sec:pseudo-labels}

This section addresses \textbf{RQ3}: it tackles the Sim2Real gap directly through
\ac{UDA} by self-training on \emph{unlabeled real} landing sequences. A
confidence-filtered self-training loop is developed and evaluated first; a
tracker-augmented extension that exploits temporal consistency is then introduced and
is \emph{ongoing}.

\subsection{Unlabeled Real Data}
\label{subsec:pl-data}
The target-domain data are the unlabeled real landing sequences described in
\Cref{subsec:setup-data}: public YouTube footage from airports and channels
\emph{disjoint} from the LARD-V1 test set, spanning the nominal, glare, rain and fog
scenarios and decomposed into approximately 1800 frames. A subset forms the custom
landing-sequence test set (glare, fog, rain) used alongside the LARD-V1 evaluation.
Keeping the sources disjoint from the test set avoids leakage and tests genuine
generalization.

\subsection{Confidence-Filtered Self-Training Loop}
\label{subsec:pl-selftraining}
The loop (\Cref{fig:pl-loop}) runs for \textbf{4 iterations}; each iteration
fine-tunes a fresh YOLOv8-pose model for 200 epochs.
\begin{enumerate}
    \item \textbf{Iteration 0} trains on the synthetic \texttt{combined} dataset only
    (equal numbers of images from the five sources).
    \item The trained model predicts runway pose on the unlabeled real frames, each
    prediction carrying a confidence score.
    \item Predictions with confidence below a threshold $\tau$ are discarded; the
    retained $m$ (image, pseudo-label) pairs are moved into the training set and an
    equal number of synthetic images is removed, keeping the training size fixed at
    5000 ($5000-m$ synthetic $+\,m$ pseudo-labeled).
    \item The model is retrained on the updated set and the process repeats.
\end{enumerate}
Two thresholds are compared, \textbf{$\tau=0.65$} and \textbf{$\tau=0.8$}, trading
pseudo-label \emph{quantity} against \emph{quality}.

\begin{figure}[th]
    \centering
    \plotslot{Self-training loop: LARD-V2 synthetic pool $\to$ train-set (5k) $\to$
    YOLOv8-pose $\to$ predictions+confidence on unlabeled YouTube frames $\to$
    confidence$>\tau$ filter $\to$ pseudo-labeled data back into train-set; evaluated
    on the LARD test set and the custom landing-sequence test set.}
    \caption{Confidence-filtered self-training pipeline (4 iterations, thresholds
    $\tau\in\{0.65,0.8\}$).}
    \label{fig:pl-loop}
\end{figure}

\subsection{Results on the LARD-V1 Test Set}
\label{subsec:pl-lard}
The two thresholds yield similar results on the LARD test set. On the nominal split
the gain over the iteration-0 baseline is small, but on the adverse splits
self-training improves mAP@[.50:.95] by more than \textbf{10 points}, with most of the
improvement realized already at the first iteration (\Cref{fig:pl-lard}). Self-training
on unlabeled real data thus transfers real-domain appearance far more effectively than
weather augmentation did in \Cref{sec:weather-aug}.

\begin{figure}[th]
    \centering
    \plotslot{mAP@[.50:.95] by iteration (0--3) for the LARD test set, per scenario
    (nominal, fog, rain, glare, night), thresholds $0.65$ vs.\ $0.8$.}
    \caption{Self-training on the LARD test set: large gains on adverse splits,
    marginal on nominal; thresholds behave similarly.}
    \label{fig:pl-lard}
\end{figure}

\subsection{Results on the Custom Landing-Sequence Test Set}
\label{subsec:pl-custom}
On the custom landing-sequence test set the gains are considerably larger, and the
threshold matters more --- particularly for \emph{fog}. Overall the baseline is
exceeded far more than on the LARD test set, by up to \textbf{35 mAP points} on the
glare split (\Cref{fig:pl-custom}). This is expected: the custom sequences come from
the same distribution as the pseudo-labels, so self-training adapts to them most
strongly.

\begin{figure}[th]
    \centering
    \plotslot{mAP@[.50:.95] by iteration for the custom test set (fog, rain, glare),
    thresholds $0.65$ vs.\ $0.8$; up to $+35$ mAP on glare.}
    \caption{Self-training on the custom landing-sequence test set: larger gains and
    stronger threshold dependence than on LARD.}
    \label{fig:pl-custom}
\end{figure}

\subsection{Pseudo-Label Composition}
\label{subsec:pl-composition}
Analyzing which pseudo-labels enter the training set (at $\tau=0.8$) reveals a
subtlety of self-training. The \emph{rain} pseudo-labels are strongly represented at
iterations~1 and~2 (both in training-set proportion and acceptance ratio), yet the
best rain-split performance is a modest $+7$ mAP reached already by the
iteration-1 model. In other words, the iteration-0 predictions provided the most
\emph{useful} rain pseudo-labels, even though later iterations accumulate many more of
them --- a reminder that pseudo-label \emph{quantity} does not equal quality and that
self-training is exposed to confirmation bias.

\plotslot{Per-scenario pseudo-label train-set proportion and acceptance ratio by
iteration ($\tau=0.8$): rain dominates iterations 1--2.}

\subsection{Tracker-Augmented Temporal Pseudo-Labels (Ongoing)}
\label{subsec:pl-tracker}
So far the pseudo-labels use the per-frame YOLOv8-pose predictions without any
additional filtering. Because the data are \emph{video} landing sequences, combining
the detector with an object tracker can exploit \emph{temporal consistency}: a
high-confidence detection can be propagated \emph{forward and backward} through the
sequence to recover the runway in frames where the detector alone is uncertain (e.g.\
distant or transiently degraded). This can be interpreted as distilling a pre-trained
tracker into temporally consistent 4-corner pseudo-labels. A first qualitative
comparison of YOLOv8-pose against YOLOv8-pose~+~OSTracker~\cite{ostrack2022} shows the
tracker recovering runways that the single-frame detector misses
(\Cref{fig:pl-tracker}).

\begin{figure}[th]
    \centering
    \plotslot{Qualitative comparison: YOLOv8-pose vs.\ YOLOv8-pose~+~OSTracker on
    landing-sequence frames; the tracker recovers the runway in low-visibility
    frames.}
    \caption{Tracker-augmented pseudo-labeling (ongoing): temporal propagation
    recovers detections missed per-frame.}
    \label{fig:pl-tracker}
\end{figure}

This extension is \emph{in progress}. The plan is to regenerate the pseudo-labels with
tracker-based temporal consistency and repeat the self-training loop for both
thresholds, comparing against the confidence-only variant above.
\todo[inline]{Once complete: report tracker-augmented self-training results per
scenario for both thresholds, and quantify the benefit of temporal consistency,
especially for distant runways.}

\subsection{Discussion}
\label{subsec:pl-discussion}
Self-training uses real data directly and therefore attacks the domain shift more
effectively than augmentation, delivering double-digit mAP gains on the adverse
splits and up to $+35$ mAP on in-distribution landing sequences, at the cost of
sensitivity to pseudo-label noise (as the rain-composition analysis shows). The
threshold choice is nearly neutral on LARD but matters on the custom set. The
tracker-augmented variant is expected to improve pseudo-label completeness in exactly
the hard, distant-runway regime identified in \Cref{subsec:gap-distance,subsec:gap-quality}.
